% Generated draft: requires mathematical and source review.
\documentclass[11pt]{article}
\usepackage[margin=25mm]{geometry}
\usepackage{fontspec}
\setmainfont{Latin Modern Roman}
\usepackage{amsmath,amssymb,mathtools,unicode-math}
\setmathfont{Latin Modern Math}
\usepackage{enumitem,needspace,xurl,hyperref}
\setlength{\parindent}{0pt}
\setlength{\parskip}{5pt}
\setlength{\emergencystretch}{3em}
\newcommand{\N}{\mathbb{N}}
\newcommand{\Z}{\mathbb{Z}}
\newcommand{\Q}{\mathbb{Q}}
\newcommand{\R}{\mathbb{R}}
\newcommand{\C}{\mathbb{C}}
\newcommand{\F}{\mathbb{F}}
\newcommand{\A}{\mathbb{A}}
\newcommand{\PP}{\mathbb{P}}
\newcommand{\E}{\mathbb{E}}
\newcommand{\eps}{\varepsilon}
\newcommand{\dd}{\,\mathrm{d}}
\begin{document}
\section*{UnsolvedMath ID: 2684}
\textbf{Source problem code:} catalogue entry
\subsection{Definitions and mathematical statement}
Let $K\subset S^3$ be a knot and $X_K$ its exterior, with boundary a torus. Let $\mathcal S_q(\partial X_K)$ denote the Kauffman bracket skein algebra of that torus, with the quantum parameter convention fixed consistently throughout. Its product comes from stacking skeins. The boundary algebra acts on the skein module of the exterior. The noncommutative $A$-ideal is the left ideal
\[
 I_K=\{P\in\mathcal S_q(\partial X_K):P\cdot\varnothing=0
 \text{ in the skein module of }X_K\},
\]
where $\varnothing$ is the empty skein.

Let $J_K=(J_n(K;q))_n$ be the infinite colored Jones vector. Once its normalization, indexing, framing and the compatible boundary-algebra action on sequences are specified, define
\[
 \operatorname{Ann}(J_K)=\{P:P\cdot J_K=0\}.
\]
The algebra in this definition must be the same as, or explicitly identified with, the one used for $I_K$; passing to a localized recurrence ring is an additional operation, not an automatic equality of ideals.

\textbf{Part (a), imported ideal-equality conjecture.} For every knot $K$, is $I_K$ exactly the annihilator of the colored Jones vector in this compatible algebraic setting?
\[
 I_K=\operatorname{Ann}(J_K).
\]
The background identifies a version for $\mathfrak{sl}_2$ in Frohman--Gelca--Lofaro and a full version in Sikora, but the short statement does not specify a more general group/Lie-algebra parameter. No such parameterized generalization is reconstructed here.

For part (b), choose meridian and longitude coordinates, denoted $M,L$. In the standard eigenvalue description, the classical $A$-polynomial records the one-dimensional boundary eigenvalue locus arising from $\operatorname{SL}_2(\C)$ representations of $\pi_1(X_K)$: boundary eigenvalues of the meridian and longitude are $M,L$. The convention for components, possible abelian factor, monomial units and coordinate choices must be fixed; the source's character-variety description does not supply all these choices.

An ordered recurrence polynomial can be written
\[
 \alpha(K,q;M,L)=\sum_{j=0}^{d}a_j(q,M)L^j,
 \qquad\alpha\cdot J_K=0,
\]
where $L$ shifts the color index and $M$ acts by the corresponding color-dependent quantum monomial in a chosen convention. The source requires nonzero $\alpha$ of minimal $L$-degree with coprime coefficients, but does not specify the coefficient ring, localization, quantum commutation relation or uniqueness normalization. These are material: one must not choose a relation incompatible with commutative specialization at $q=-1$.

\textbf{Part (b), AJ conjecture (Garoufalidis).} With matching normalizations and a defined specialization, does
\[
 \alpha(K,-1;M,L)=A(K;L,M)
 \quad\text{up to a nonzero factor depending only on }M
\]
hold for every knot? The order of arguments on the two sides and the specialization at $-1$ are retained exactly. The allowed coefficient domain for the $M$-factor is not stated; the often-used interpretation as equality up to a nonzero element of $\C(M)$ is a convention requiring confirmation, not an exact polynomial equality silently imposed here. Specialization, denominator clearing, primitive normalization and any component factors must all be controlled.

The background reports that Frohman--Gelca--Lofaro relate the ideal to colored Jones recurrences, computations of $I_K$ for the trefoil $3_1$, $(2,2p+1)$ torus knots and figure-eight $4_1$, and general nontriviality of the noncommutative ideal as unresolved. The cited $q$-holonomicity of the colored Jones function supplies recurrence information in its appropriate operator setting; it is not silently identified with nontriviality or equality in an unspecified unlocalized skein algebra.

The literature triage dated 2026-08-17 reports AJ results for specified 2-bridge and pretzel families under the cited papers' hypotheses, and a diagrammatic framework rather than a universal proof. The background also cites a physics interpretation. Exact family conditions are not reproduced and are not generalized to all such knots. The review leaves general ideal equality/nontriviality and AJ specialization open, with possible $M$-factors explicitly unresolved. Its partially-solved label and the registry's open label are supplied dated assessments. No proof, PDF or source was independently inspected.
\subsection{Short English statement}
For each knot, compare the boundary skein operators killing the empty skein with the operators annihilating its colored Jones sequence. The AJ part asks whether a minimal quantum recurrence specializes at $q=-1$ to the classical $A$-polynomial, up to a factor in $M$. The source leaves important ring, normalization and component conventions unresolved and reports results only for specified knot families.
\subsection{Sources}
\begin{itemize}
\item UnsolvedMath local registry, data/unsolvedmath/problems.json, permanent ID 2684, source code catalogue entry, ``Kirby Problem 1.25,'' including both conjectures, all literature notes and triage dated 2026-08-17.
\item \emph{K3: A New Problem List in Low-Dimensional Topology}, Problem 1.25 (2026), supplied link \url{https://math.berkeley.edu/sites/default/files/surv-295-ruberman-watermarked-author-pdf.pdf}. The review cites the distinction between general conjectures and family results; the PDF was not opened.
\item [FGL02] Charles Frohman, Răzvan Gelca and Walter Lofaro, ``The A-polynomial from the noncommutative viewpoint,'' \emph{Trans. Amer. Math. Soc.} 354(2) (2002), 735--747, DOI 10.1090/S0002-9947-01-02889-6, supplied review link \url{https://doi.org/10.1090/S0002-9947-01-02889-6}.
\item [Sik08] Adam S. Sikora, ``Quantizations of character varieties and quantum knot invariants'' (2008), arXiv:0807.0943.
\item [Gar04] Stavros Garoufalidis, ``On the characteristic and deformation varieties of a knot,'' \emph{Proceedings of the Casson Fest}, Geom. Topol. Monogr. 7 (2004), 291--309, DOI 10.2140/gtm.2004.7.291.
\item [Gel02] Răzvan Gelca, ``Non-commutative trigonometry and the A-polynomial of the trefoil knot,'' \emph{Math. Proc. Cambridge Philos. Soc.} 133(2) (2002), 311--323, DOI 10.1017/S0305004102006047.
\item [GS03] Răzvan Gelca and Jeremy Sain, ``The noncommutative A-ideal of a $(2,2p+1)$-torus knot determines its Jones polynomial,'' \emph{J. Knot Theory Ramifications} 12(2) (2003), 187--201, DOI 10.1142/S021821650300238X.
\item [GS04] Răzvan Gelca and Jeremy Sain, ``The computation of the non-commutative generalization of the A-polynomial of the figure-eight knot,'' \emph{J. Knot Theory Ramifications} 13(6) (2004), 785--808, DOI 10.1142/S0218216504003482.
\item [CL98] D. Cooper and D. D. Long, ``Representation theory and the A-polynomial of a knot,'' \emph{Chaos Solitons Fractals} 9(4--5) (1998), 749--763, DOI 10.1016/S0960-0779(97)00102-1.
\item [GL05] Stavros Garoufalidis and Thang T. Q. Lê, ``The colored Jones function is qholonomic,'' \emph{Geom. Topol.} 9 (2005), 1253--1293, DOI 10.2140/gt.2005.9.1253.
\item [LZ17] Thang T. Q. Lê and Xingru Zhang, ``Character varieties, A-polynomials and the AJ conjecture,'' \emph{Algebr. Geom. Topol.} 17(1) (2017), 157--188, DOI 10.2140/agt.2017.17.157, supplied review link \url{https://doi.org/10.2140/agt.2017.17.157}.
\item [LT15] Thang T. Q. Le and Anh T. Tran, ``On the AJ conjecture for knots,'' \emph{Indiana Univ. Math. J.} 64(4) (2015), 1103--1151, with an appendix jointly with Vu Q. Huynh, DOI 10.1512/iumj.2015.64.5602, supplied review link \url{https://doi.org/10.1512/iumj.2015.64.5602}.
\item [DG20] Renaud Detcherry and Stavros Garoufalidis, ``A diagrammatic approach to the AJ conjecture,'' \emph{Math. Ann.} 378(1--2) (2020), 447--484, DOI 10.1007/s00208-020-02028-y.
\item [Guk05] Sergei Gukov, ``Three-dimensional quantum gravity, Chern--Simons theory, and the A-polynomial,'' \emph{Comm. Math. Phys.} 255(3) (2005), 577--627, DOI 10.1007/s00220-005-1312-y. The additional inline label [Lê06] has no bibliographic entry supplied.
\item Ulam AI and the UnsolvedMath Contributors, \emph{UnsolvedMath: A Curated Collection of Open Mathematics Problems}. Dataset curation and original metadata are supplied under CC BY 4.0; referenced source material retains its own terms. No external reference was independently consulted for this draft.
\end{itemize}
\end{document}
